---
format: yohaku/0.2
target: text
macro: '[[macros/raw]]'
---

## 正文
````text
"""Request-owned stdlib observation; no global patches or payload logging."""
import errno
import json
import threading
import socket
import urllib.error
import hashlib
import http.client
import os
import ssl
import time
import urllib.parse
import urllib.request

class TransportTrace:
    def __init__(self, url, timeout=None, request_bytes=None, messages=None, recovery=None):
        self.recovery=recovery
        self.started=time.monotonic();self.events=[];self.dropped=0
        self.counts={};self.failed_phase=None;self.origin_status=None
        self.target=urllib.parse.urlsplit(url).hostname
        self.metadata={'timeout_seconds':timeout,'request_bytes':request_bytes,
            'connection_policy':'stdlib urllib new connection per open; Connection: close',
            'remote_receipt':'unknown; local send completion does not prove server receipt'}
        proxies=urllib.request.getproxies()
        self.metadata['proxy_schemes']=sorted(k for k in proxies if k in ('http','https','all','no'))
        self.metadata['target_bypass']=bool(urllib.request.proxy_bypass(self.target or ''))
        if recovery:self.metadata['recovery']={'attempt':getattr(recovery,'attempt',1),'max_attempts':recovery.max_retries+1,'previous_failures':list(recovery.request_events)}
        self.metadata['ca_environment_set']={k:bool(os.environ.get(k)) for k in ('SSL_CERT_FILE','SSL_CERT_DIR')}
        self.metadata['default_ca_paths']={'cafile':ssl.get_default_verify_paths().cafile,'capath':ssl.get_default_verify_paths().capath}
        if isinstance(messages,list):
            counts={role:sum(isinstance(m,dict) and m.get('role')==role for m in messages) for role in ('user','assistant','tool','system')}
            self.metadata['message_role_counts']=counts
            confirmed=[]
            for message in messages:
                if not isinstance(message,dict):continue
                if message.get('role')=='tool':confirmed.append((message.get('tool_call_id'),message.get('content')))
                elif message.get('type')=='function_call_output':confirmed.append((message.get('call_id'),message.get('output')))
                content=message.get('content')
                if message.get('role')=='user' and isinstance(content,list):
                    confirmed.extend((block.get('tool_use_id'),block.get('content')) for block in content
                        if isinstance(block,dict) and block.get('type')=='tool_result')
            self.metadata['confirmed_tool_results']=len(confirmed)
            if confirmed:
                for name,value in zip(('tool_call_id','content'),confirmed[-1]):
                    if isinstance(value,str):self.metadata['last_tool_'+name+'_sha256']=hashlib.sha256(value.encode()).hexdigest()
    def event(self, phase, **fields):
        self.counts[phase]=self.counts.get(phase,0)+1
        entry={'phase':phase,'elapsed_seconds':round(time.monotonic()-self.started,6),**fields}
        if len(self.events)<64:self.events.append(entry)
        else:self.dropped+=1
    def failure(self, phase, error):
        if self.failed_phase is None:self.failed_phase=phase
        self.event(phase+'_failed',exception_type=type(error).__name__)
    def snapshot(self):
        return {'metadata':self.metadata,'events':self.events,'events_dropped':self.dropped,
            'counts':self.counts,'failed_phase':self.failed_phase,'origin_http_status':self.origin_status,
            'elapsed_seconds':round(time.monotonic()-self.started,6)}

class _Context:
    def __init__(self, original, trace):self.original=original;self.trace=trace
    def __getattr__(self,name):return getattr(self.original,name)
    def wrap_socket(self,*args,**kwargs):
        t=self.trace;t.event('tls_wrap_handshake_start',verify_mode=int(self.original.verify_mode),check_hostname=self.original.check_hostname,
            handshake_on_connect=kwargs.get('do_handshake_on_connect',True))
        try:sock=self.original.wrap_socket(*args,**kwargs)
        except BaseException as error:t.failure('tls_wrap_handshake',error);raise
        if t.recovery:t.recovery.track(sock)
        t.event('tls_wrap_handshake_complete',tls_version=sock.version());return sock

class _Connection:
    def __init__(self, trace, *args, **kwargs):
        self.trace=trace;self.in_tunnel=False
        super().__init__(*args,**kwargs)
        original=self._create_connection
        def connect(*args,**kwargs):
            trace.event('dns_tcp_connect_start')
            try:sock=original(*args,**kwargs)
            except BaseException as error:trace.failure('dns_tcp_connect',error);raise
            if trace.recovery:trace.recovery.track(sock)
            trace.event('dns_tcp_connect_complete');return sock
        self._create_connection=connect
        if hasattr(self,'_context'):self._context=_Context(self._context,trace)
        trace.event('connection_created',target_host_matches=self.host==trace.target,port=self.port)
    def connect(self):
        self.trace.event('connection_start',proxy_tunnel=bool(self._tunnel_host))
        return super().connect()
    def _tunnel(self):
        self.in_tunnel=True;self.trace.event('proxy_tunnel_start')
        try:result=super()._tunnel()
        except BaseException as error:self.trace.failure('proxy_tunnel',error);raise
        else:self.trace.event('proxy_tunnel_complete');return result
        finally:self.in_tunnel=False
    def send(self,data):
        t=self.trace;phase='proxy_send' if self.in_tunnel else 'request_send'
        t.event(phase+'_start',socket_present=self.sock is not None)
        try:result=super().send(data)
        except BaseException as error:t.failure(phase,error);raise
        t.event(phase+'_complete');return result
    def getresponse(self):
        t=self.trace;t.event('response_headers_start')
        try:response=super().getresponse()
        except BaseException as error:t.failure('response_headers',error);raise
        t.origin_status=response.status;t.event('response_headers_complete',status=response.status)
        return response

class ObservedHTTP(_Connection,http.client.HTTPConnection):pass
class ObservedHTTPS(_Connection,http.client.HTTPSConnection):pass

class HTTPHandler(urllib.request.HTTPHandler):
    def __init__(self,trace):super().__init__();self.trace=trace
    def http_open(self,req):return self.do_open(lambda *a,**k:ObservedHTTP(self.trace,*a,**k),req)

class HTTPSHandler(urllib.request.HTTPSHandler):
    def __init__(self,trace,**kwargs):super().__init__(**kwargs);self.trace=trace
    def https_open(self,req):return self.do_open(lambda *a,**k:ObservedHTTPS(self.trace,*a,**k),req,context=self._context)

def handlers(trace):return HTTPHandler(trace),HTTPSHandler(trace)


class RecoveryCancelled(Exception):pass


TURN_TIMEOUT_DEFAULT = None
TURN_TIMEOUT_LONG = 86400

def turn_seconds(config):
    # Missing/null means unlimited. A legacy long preset remains explicit opt-in.
    if 'turn_timeout_long' in config and type(config['turn_timeout_long']) is not bool:
        raise ValueError('turn_timeout_long must be a boolean')
    if 'turn_timeout' in config and config['turn_timeout'] is None:return None
    seconds=config.get('turn_timeout')
    if seconds is None:return TURN_TIMEOUT_LONG if config.get('turn_timeout_long') is True else None
    if type(seconds) is not int or seconds<=0:
        raise ValueError('turn_timeout must be a positive integer number of seconds or null for unlimited')
    return max(seconds,TURN_TIMEOUT_LONG) if config.get('turn_timeout_long') is True else seconds

def turn_policy(config):
    seconds=turn_seconds(config)
    key='turn_timeout' if 'turn_timeout' in config else ('turn_timeout_long' if config.get('turn_timeout_long') is True else 'default')
    provenance=config.get('_turn_budget_provenance',{})
    source=provenance.get(key,'wire.initial.config.'+key if key!='default' else 'worker default: unlimited')
    return {'mode':'unlimited' if seconds is None else 'limited','seconds':seconds,'source':source,
            'configuration_reads':provenance.get('configuration_reads',[])}

def resolve_turn_config(initial,project,host_config=None):
    """Refresh only public total-limit keys on native startup, even with an old host.

    Model routing, credentials, permissions, request and command timeouts are untouched.
    Standalone/test callers without a real native world retain their supplied config.
    """
    from pathlib import Path
    config=dict(initial['config']);context=initial.get('context',{})
    root=context.get('root') if isinstance(context,dict) else None
    if not isinstance(root,str) or not Path(root).is_absolute() or not Path(root).is_dir():return config
    host=Path(host_config) if host_config else Path(project)/'config/ai.json';world=Path(root)/'memory/socius.json'
    if not host.is_file():return config
    keys=('turn_timeout','turn_timeout_long')
    for key in keys:config.pop(key,None)
    sources={'configuration_reads':[]}
    for path,is_world in ((host,False),(world,True)):
        if not path.is_file():continue
        value=json.loads(path.read_text())
        if not isinstance(value,dict):raise ValueError('Invalid limit configuration: '+str(path))
        selected=value.get('chat',{}).get('limits',{}) if is_world else value
        if not isinstance(selected,dict):raise ValueError('chat.limits must be an object')
        sources['configuration_reads'].append(str(path))
        for key in keys:
            if key in selected:config[key]=selected[key];sources[key]=str(path)+('#/chat/limits/' if is_world else '#/')+key
    config['_turn_budget_provenance']=sources
    turn_seconds(config) # Fail before connect, never substitute a 900s fallback.
    return config

class ReplayBlocked(Exception):pass

class ToolReplayGuard:
    """Worker-local ambiguity barrier for one recovered logical model request.

    Read-only reconciliation remains available. A repeated effect after recovery
    pauses instead of guessing intent. A fresh request without a retry, or a new
    user turn, does not permanently suppress an intentional repeated operation.
    Only fingerprints are retained here; authoritative results stay in history.
    """
    READ_ONLY=frozenset(('context','list_files','list_directory','read_document',
        'read_text','read_range','search_text','search_notes','list_entry_nodes',
        'find_nodes','get_node','walk_nodes','list_tools','tool_status','inspect_tool',
        'terminal_status','terminal_list','terminal_read','codex_model_status','codex_model_list'))
    def __init__(self):self.completed=set();self.before=set();self.protected=set()
    def begin_request(self):self.before=set(self.completed);self.protected=set()
    def arm(self):self.protected.update(self.before)
    def execute(self,name,args,operation):
        if name in self.READ_ONLY:return operation()
        fingerprint=hashlib.sha256(json.dumps([name,args],sort_keys=True,ensure_ascii=False,separators=(',',':')).encode()).hexdigest()
        if fingerprint in self.protected:
            raise ReplayBlocked('Possible repeated operation after model recovery: '+name+
                '; fingerprint='+fingerprint+'; previous dispatch has a confirmed result; no replay performed. '+
                'Reconcile saved state using read_range, tool_status or terminal_status. '+
                'If intent remains ambiguous, pause and ask the user; a new user turn permits intentional repetition.')
        # Even an ok=false result can contain a saved write or a failed process.
        result=operation();self.completed.add(fingerprint);return result


def retryable(error):
    """Only temporary transport failures, never malformed/semantic model output."""
    seen=set()
    while error is not None and id(error) not in seen:
        seen.add(id(error))
        if isinstance(error,(RecoveryCancelled,ssl.SSLCertVerificationError)):return False
        if isinstance(error,urllib.error.HTTPError):return error.code in (408,429,500,502,503,504)
        if isinstance(error,(TimeoutError,ConnectionError,http.client.IncompleteRead,
                             http.client.RemoteDisconnected,ssl.SSLEOFError)):return True
        if isinstance(error,OSError) and error.errno in (errno.EAGAIN,errno.ECONNRESET,errno.ECONNABORTED,
                errno.EPIPE,errno.ETIMEDOUT,errno.ECONNREFUSED,errno.ENETUNREACH,errno.EHOSTUNREACH):return True
        reason=getattr(error,'reason',None)
        error=reason if isinstance(reason,BaseException) else error.__cause__
    return False


class Recovery:
    """One owner per worker. Retry model I/O only; call budgets count logical rounds.

    Reuse Codex's request_max_retries spelling. Native Codex retries stay zero
    when our relay owns recovery. No retries are allowed after delivery to tools.
    """
    def __init__(self,config=None):
        retries=(config or {}).get('request_max_retries',2)
        if type(retries) is not int or not 0<=retries<=5:
            raise ValueError('request_max_retries must be an integer from 0 to 5')
        self.max_retries=retries;self.cancelled=threading.Event();self.deadline=None
        self.lock=threading.RLock();self.sockets=set();self.events=[];self.request_events=[]
        self.notify=lambda event:None
        self.tool_guard=ToolReplayGuard();self.expired=False;self.timer=None

    def set_budget(self,seconds):
        if self.timer:self.timer.cancel()
        self.deadline=None if seconds is None else time.monotonic()+seconds
        if seconds is None:return
        def expire():
            self.expired=True
            self.cancel()
        self.timer=threading.Timer(seconds,expire);self.timer.daemon=True;self.timer.start()

    def finish_budget(self):
        if self.timer:self.timer.cancel()

    def check(self):
        if self.expired:raise RecoveryCancelled('Turn time budget exhausted; completed tools remain saved; inspect running tools separately; no replay')
        if self.cancelled.is_set():raise RecoveryCancelled('Model request stopped; no retry or tool replay')
        if self.deadline is not None and time.monotonic()>=self.deadline:
            raise RecoveryCancelled('Turn time budget exhausted; saved tool results remain; no replay')

    def track(self,sock):
        with self.lock:self.sockets.add(sock)
        if self.cancelled.is_set():
            try:sock.shutdown(socket.SHUT_RDWR)
            except OSError:pass
            self.check()

    def cancel(self):
        self.cancelled.set()
        with self.lock:sockets=list(self.sockets)
        for sock in sockets:
            try:sock.shutdown(socket.SHUT_RDWR)
            except OSError:pass

    def run(self,operation,details=lambda:{},notify=None):
        self.request_events=[]
        self.check()
        self.tool_guard.begin_request()
        for attempt in range(1,self.max_retries+2):
            self.attempt=attempt
            try:
                if attempt>1:
                    (notify or self.notify)({'event':'status','text':'Model request retry '+str(attempt)+'/'+str(self.max_retries+1)+' requesting; Stop available'})
                self.check();value=operation();self.check();return value
            except Exception as error:
                self.check()
                again=retryable(error) and attempt<=self.max_retries
                event={'attempt':attempt,'max_attempts':self.max_retries+1,
                    'exception_type':type(error.__cause__ or error).__name__,
                    'will_retry':again,**details()}
                self.events.append(event);self.events=self.events[-32:]
                self.request_events.append(event)
                emit=notify or self.notify
                emit({'event':'recovery','backend':'model-request',**event})
                emit({'event':'status','text':'Model request attempt '+str(attempt)+'/'+str(self.max_retries+1)+
                    ' failed: '+event['exception_type']+'; phase='+str(event.get('phase'))+
                    ('; retry pending; Stop available' if again else '; recovery exhausted or not retryable; saved tool results remain')})
                if not again:raise
                self.tool_guard.arm()
                delay=min(.5*2**(attempt-1),2)
                if self.deadline is not None:delay=min(delay,max(0,self.deadline-time.monotonic()))
                if self.cancelled.wait(delay):self.check()
                self.check()
            finally:
                with self.lock:self.sockets.clear()
````

## 输出与关联
```yaml
output: "agent/yohaku_provider_transport.py"
```

## AI 笔记
```yaml
{"source_version": 1, "writer": "Yohaku release packaging", "summary": "Portable release source. Scope and packaging adjustments are documented in PROVENANCE.md."}
```
